\chapter*{Conclusion}
\addcontentsline{toc}{chapter}{Conclusion}
\label{ch:conclusion}

This thesis has developed a unified framework for understanding the
vacuum structure of quantum field theories and its role in the physics
of the early universe, proceeding from the abstract machinery of the
effective action to the concrete phenomenon of Hawking radiation.

We began in Part~I by constructing the effective action formalism from
first principles. Starting from the generating functional $Z[J]$ and
passing through the connected functional $W[J]$ via a Legendre
transform, we arrived at the effective action $\Gamma[\varphi_c]$,
whose equation of motion governs the full quantum dynamics of the
field. The central object for our purposes was the effective potential
$V_{\text{eff}}(\varphi_c)$, whose minima determine the true vacuum of
the quantum theory. We gave it two complementary physical
interpretations: as the generating functional for 1PI Green's functions
at zero external momenta, and as the minimum energy density over all
normalized states with a given field expectation value. The one-loop
computation in scalar field theory, first using a hard momentum cutoff
and then using dimensional regularization and the $\overline{\text{MS}}$
scheme, gave the Coleman-Weinberg potential, and we applied it to gauge
theories including the GSW electroweak model and GUT models. The
striking phenomenon of dimensional transmutation in massless scalar
electrodynamics showed how quantum corrections can generate a mass
scale and a nontrivial vacuum structure in a classically scale-invariant
theory.

Part~II extended the formalism to finite temperature using the
imaginary-time formalism of thermal field theory. The key structural
change is that the periodicity of fields in Euclidean time with period
$\beta = 1/T$, the KMS condition, discretizes the loop momenta into
Matsubara frequencies. The finite-temperature effective potential
acquires thermal corrections that generate a positive mass term
proportional to $T^2\varphi_c^2$, driving symmetry restoration at high
temperatures even in theories that are spontaneously broken at $T = 0$.
We computed these corrections explicitly for scalar field theory, the
GSW model, and minimal SU(5), and derived the critical temperature
$T_c$ at which the two phases become degenerate.

Part~III placed this analysis in its cosmological context. The thermal
history of the early universe is a sequence of symmetry breaking
transitions, each driven by the temperature-dependent effective
potential falling below its critical value. We revisited the
Coleman-Weinberg potential in the $\overline{\text{MS}}$ scheme and
introduced the thermal functions $J_B$ and $J_F$, whose
high-temperature expansions reveal the full structure of the
finite-temperature corrections including the cubic term $\sim T|m|^3$
that is responsible for generating a potential barrier and hence a
first-order phase transition. The Coleman-Weinberg model at finite
temperature provided a concrete example of how a classically
scale-invariant theory develops a preferred scale and undergoes a
first-order transition as the universe cools.

Part~IV turned to the dynamics of false vacuum decay and bubble
nucleation, the mechanism by which a first-order phase transition
actually proceeds. The universe, trapped in a metastable false vacuum,
tunnels to the true vacuum through the nucleation of bubbles governed
by the bounce solution of the Euclidean field equations. We analyzed
the $O(4)$-symmetric bounce at zero temperature and the
$O(3)$-symmetric bounce at finite temperature, computed the nucleation
rate, and characterized the transition by the parameters $\alpha$ and
$\tilde\beta$ that determine the gravitational wave spectrum produced
by bubble collisions. This provides a direct observational window into
the thermal history of the early universe, potentially accessible to
future detectors such as LISA.

Part~V addressed the fundamental limitation of everything that preceded
it: the assumption that the background spacetime is fixed and flat. In
the early universe and near compact objects, the curvature of spacetime
is not negligible, and quantum field theory must be formulated on a
curved background. We developed the semiclassical gravity framework, in
which the geometry is treated classically via the Einstein equations
but matter fields are quantized. The central new feature is the
observer-dependence of the vacuum state: in the absence of a global
Poincar\'{e} symmetry, different observers define inequivalent mode
decompositions related by Bogoliubov transformations, and the vacuum
of one observer contains particles as seen by another. We made this
precise through the moving mirror model and the Unruh effect, before
deriving the Hawking effect: a black hole formed by gravitational
collapse emits a steady thermal flux of radiation at temperature
$T_H = \kappa/2\pi$, where $\kappa$ is the surface gravity. The
Bekenstein-Hawking entropy $S_{BH} = A/4$ and the generalized second
law complete the thermodynamic picture of black holes.

Several threads connect these results into a unified story. The
effective potential, computed perturbatively in flat spacetime, acquires
curvature corrections in a curved background through the non-minimal
coupling $\xi R\phi^2$, which can shift the phase structure of the
theory. The thermal functions $J_B$ and $J_F$ that appear in the
finite-temperature effective potential are the same integrals that
appear in the Stefan-Boltzmann free energy of the Hawking radiation.
The Hawking temperature $T_H = 1/8\pi M$ satisfies $T = dM/dS_{BH}$,
confirming the thermodynamic consistency of the Bekenstein-Hawking
entropy. And the same mechanism of particle production near a causal
horizon that gives rise to Hawking radiation also gives rise to the
Gibbons-Hawking temperature $T = H/2\pi$ of de Sitter space, which
seeds the primordial density perturbations generated during inflation.

The work of this thesis has been primarily foundational, building the
theoretical tools needed to address these questions rather than
pursuing any single one of them to its conclusion. The effective action
formalism, the imaginary-time thermal field theory, and the Bogoliubov
transformation framework are the basic ingredients of any serious
treatment of quantum fields in the early universe. Having developed
them systematically and seen how they connect to one another, the
natural next step is to apply them more carefully to specific physical
scenarios, refining the approximations made here and confronting the
results with observation.
